\label{sec:antihydra}

In this section, we define the Antihydra Turing machine and prove that its halting behavior is equivalent to a specific condition on the iterates of a mathematical map $A: \mathbb{N} \times \mathbb{N} \to \mathbb{N} \times \mathbb{N}$.

\subsection{The Antihydra Turing Machine}

The Antihydra machine is a 6-state, 2-symbol Turing machine discovered by Shawn Ligocki. It is defined by the following transition table.

\begin{definition}
  \label{def:transition}
  \lean{Antihydra.transition}
  \leanok
  The state set is $\{A, B, C, D, E, F\}$ and the tape alphabet is $\{0, 1\}$. The transition function $\delta: Q \times \Sigma \to (Q \cup \{HALT\}) \times \Sigma \times \{L, R\}$ is given by:
  \begin{center}
  \begin{tabular}{|c|c||c|c|c|}
    \hline
    State & Symbol & Next State & Write & Move \\
    \hline
    A & 0 & B & 1 & R \\
    A & 1 & A & 1 & R \\
    B & 0 & C & 0 & L \\
    B & 1 & E & 1 & L \\
    C & 0 & D & 1 & L \\
    C & 1 & C & 1 & L \\
    D & 0 & A & 1 & L \\
    D & 1 & B & 0 & L \\
    E & 0 & F & 1 & L \\
    E & 1 & E & 1 & R \\
    F & 0 & HALT & 0 & R \\
    F & 1 & A & 0 & R \\
    \hline
  \end{tabular}
  \end{center}
\end{definition}

\begin{definition}
  \label{def:Config}
  \lean{Antihydra.Config}
  \leanok
  A configuration of the machine consists of the current state, the contents of the tape to the left and right of the head, and the symbol currently under the head.
\end{definition}

\subsection{The Mathematical Model}

The high-level behavior of the Antihydra machine can be modeled as a dynamical system on pairs of natural numbers.

\begin{definition}
  \label{def:A_map}
  \lean{Antihydra.A}
  \leanok
  The map $A: \mathbb{N} \times \mathbb{N} \to \mathbb{N} \times \mathbb{N}$ is defined for $(a, b) \in \mathbb{N}^2$ by:
  \[
  A(a, b) = \begin{cases}
    (3 \lfloor a/2 \rfloor + 2, b + 2) & \text{if } a \text{ is even}, \\
    (3 \lfloor a/2 \rfloor + 3, b - 1) & \text{if } a \text{ is odd}.
  \end{cases}
  \]
\end{definition}

\begin{definition}
  \label{def:mathHaltingCondition}
  \lean{Antihydra.mathHaltingCondition}
  \leanok
  The mathematical system halts if it reaches a state $(a, b)$ such that $a$ is odd, $a \ge 3$, and $b = 0$.
\end{definition}

\subsection{Equivalence of Halting}

The core result is that the Turing machine simulates the mathematical map $A$. This is proven by identifying a family of "clean" configurations that correspond to the state $(a, b)$.

\begin{definition}
  \label{def:P_Config_Pad}
  \lean{Antihydra.P_Config_Pad}
  \leanok
  For $a, b, n, p \in \mathbb{N}$, the configuration $P(b, a, n, p)$ corresponds to state $E$, reading symbol 0, with the left tape containing $1^a 0 1^b$ and the right tape containing $1^n 0^p$.
\end{definition}

\begin{theorem}
  \label{thm:tm_simulates_math}
  \lean{Antihydra.tm_simulates_math}
  \leanok
  If the machine is in a "clean" loop configuration corresponding to $(a, b)$, then after some number of steps $k > 0$, it will either reach the HALT state (if the math model halts) or return to a "clean" loop configuration corresponding to $A(a, b)$.
  \begin{proof}
    \leanok
    The proof proceeds by exhaustive analysis of the machine's execution paths (macro-steps). We show that for even $a$, the machine increases $b$ and performs a $(3/2)$ growth on $a$ (with an offset). For odd $a$, it decreases $b$ and performs a $(3/2)$ growth on $a$, unless $b=0$, in which case it transitions to state $F$ and halts.
  \end{proof}
\end{theorem}

\begin{theorem}
  \label{thm:antihydra_halts_iff}
  \lean{Antihydra.antihydra_halts_iff}
  \leanok
  The Antihydra Turing machine starting from a blank tape halts if and only if there exists some iteration $i$ such that $A^i(8, 2)$ satisfies the mathematical halting condition.
  \begin{proof}
    \leanok
    Explicit simulation for 58 steps shows that the machine reaches the configuration $P(2, 8, 0, 0)$, which corresponds to the mathematical state $(8, 2)$. By Theorem~\ref{thm:tm_simulates_math}, the machine then follows the orbit of $(8, 2)$ under $A$ until it satisfies the halting condition or continues indefinitely.
  \end{proof}
\end{theorem}
